
\section{\BookRecordIndex{T0283}求值映射与\BookRecordIndex{T0337}伪循环}\label{sec:1-3}
\zhindex{qiuzhiyingshe}{求值映射}
\zhindex{weixunhuan}{伪循环}

\BookRecordIndex{T0143}格罗莫夫--威滕不变量由\BookRecordIndex{T0283}求值映射
\[
  \mathcal M^*(A;J)\times\mathbb{CP}^1\longrightarrow M,
  \qquad (u,z)\longmapsto u(z)
\]
构造。注意，此映射通过\BookRecordIndex{T0413}重参数化群 $G$ 的作用分解，其中
\[
  \phi\cdot(u,z)=\bigl(u\circ\phi^{-1},\phi(z)\bigr).
\]
因此得到定义在商上的映射
\[
  \operatorname{ev}=\operatorname{ev}_J:
  \mathcal M^*(A;J)\times_G\mathbb{CP}^1\longrightarrow M.
\]

\begin{namedexample}
设 $V$ 是维数为 $2n-2$ 的闭\BookRecordIndex{T0358}辛流形，并令
\[
  M:=\mathbb{CP}^1\times V
\]
带有乘积\BookRecordIndex{T0363}辛形式。假设 $\pi_2(V)=0$。于是
$A:=[\mathbb{CP}^1\times\mathrm{pt}]$ 生成 $M$ 中球面 $2$-类组成的群，所以
$\omega(A)$ 必然是 $\omega$ 在\BookRecordIndex{T0285}球面类上所取的最小值。定理 A 和 B 因而蕴含：对一般的
$J$，空间 $\mathcal M^*(A;J)\times_G\mathbb{CP}^1$ 是紧流形。因为 $c_1(A)=2$，
在这一情形下，$\mathcal M^*(A;J)\times_G\mathbb{CP}^1$ 的维数为 $2n$，恰好等于
$M$ 的维数。此外，第 3 章将说明（见定理 3.1.8），不同的 $J$ 给出彼此配边的映射
$\operatorname{ev}_J$。配边映射具有相同的度，因此 $\operatorname{ev}_J$ 的度与所有
选择无关。现在，若 $J=i\times J'$ 是乘积，其中 $i$ 表示 $\mathbb{CP}^1$ 上的标准
复结构，则容易看出 $\mathcal M^*(A;J)$ 的元素具有形式
\[
  u(z)=\bigl(\phi(z),v_0\bigr),
\]
其中 $v_0\in V$ 且 $\phi\in G$。所以对这一 $J$，映射 $\operatorname{ev}_J$ 的度为
$1$；于是对每个正则的 $J$，其度都为 $1$。
\end{namedexample}

\personlater{geluomofu}{格罗莫夫}{Gromov}在 \srcite{160} 中利用这一事实证明了著名的
\BookRecordIndex{T0120}非挤压定理。

\begin{namedtheorem}
设 $V$ 是维数为 $2n-2$ 的闭\BookRecordIndex{T0358}辛流形且 $\pi_2(V)=0$。若 $\psi$ 是从半径为 $r$ 的球
$B^{2n}(r)$ 到柱体 $B^2(\lambda)\times V$ 的辛嵌入，则 $r\le\lambda$。
\end{namedtheorem}

\begin{proof}[证明略]
把圆盘 $B^2(\lambda)$ 嵌入一个面积为 $\pi\lambda^2+\varepsilon$ 的 $2$-球面
$\mathbb{CP}^1$，并令 $\omega$ 为 $\mathbb{CP}^1\times V$ 上的乘积辛结构。取
$\mathbb{CP}^1\times V$ 上一个被 $\omega$ \BookRecordIndex{T0372}驯服的\BookRecordIndex{T0088}殆复结构 $J'$，使它在 $\psi$ 的像上
等于球 $B^{2n}(r)$ 的标准结构 $J_0$ 在 $\psi$ 下的前推。由于\BookRecordIndex{T0283}求值映射
$\operatorname{ev}_{J'}$ 的度为 $1$，$\mathbb{CP}^1\times V$ 的每一点都通过一条
$J'$-全纯曲线。特别地，有一条这样的曲线，记作 $C'$，通过球心像 $\psi(0)$。经由
$\psi$ 拉回，这条曲线在球 $B^{2n}(r)$ 内给出一条通过球心的 $J_0$-全纯曲线 $C$。
因为 $J_0$ 是标准结构，$C$ 在通常意义下全纯，因而是 $B^{2n}(r)$ 中的\BookRecordIndex{T0175}极小曲面。
\BookRecordIndex{T0175}极小曲面理论中的一个标准结果断言，欧氏球中通过球心的恰当嵌入曲面，其面积至少等于
平圆盘的面积，即 $\pi r^2$。又因为 $C$ 全纯，容易验证其面积恰等于标准\BookRecordIndex{T0363}辛形式
$\omega_0$ 在它上面的积分。因此
\[
  \pi r^2
  \le \int_C\omega_0
  =\int_{\psi^{-1}(C')}\psi^*\omega
  <\int_{C'}\omega
  =\omega(A)=\pi\lambda^2+\varepsilon.
\]
中间的不等式成立，是因为 $\psi(C)$ 只是 $C'$ 的一部分。由于这对每个
$\varepsilon>0$ 都成立，结论随即得到。
\end{proof}

上述论证的更多细节见 \srcite{160,320}。第~\ref{sec:9-3} 节将完整给出一个略有不同的证明，以
辛吹起代替对\BookRecordIndex{T0175}极小曲面理论的引用。第~\ref{ch:applications-symplectic-topology} 章还完整证明格罗莫夫论文 \srcite{160} 中陈述的
另一些基础结果。例如，第~\ref{sec:9-2} 节使用第 3 章发展的 $J$-全纯圆盘理论，证明当 $n\ge2$
时，$\mathbb R^{2n}$ 中不存在闭的恰当\BookRecordIndex{T0219}拉格朗日子流形。

\paragraph{\BookRecordIndex{T0145}格罗莫夫--威滕伪循环。}
在多个点对映射 $u$ 求值往往很有用。若求值点不超过三个，则由于 $G$ 在球面上三重
传递，选择哪些点并不重要。不过，若求值点超过三个（或者定义域不是球面），点的选择就会
造成差别。我们把这些点纳入\BookRecordIndex{T0246}模空间本身，于是基本研究对象成为元组
$(u,z_1,\dots,z_k)$：其中 $u\in\mathcal M(A;J)$，而 $z_i\in S^2$ 是 $k$ 个两两
不同的\BookRecordIndex{T0066}标记点。所有这类元组组成的空间记为 $\widetilde{\mathcal M}_{0,k}(A;J)$，它对
$G$ 的商记为 $\mathcal M_{0,k}(A;J)$。
\footnote{这一记号的下标 $0$ 表示定义域的亏格。本书只讨论亏格为 $0$ 的情形，不过，
当 $(\Sigma,j)$ 具有任意亏格而复结构 $j$ 固定时，本书许多论述几乎逐字适用。}
与前面一样，$\mathcal M^*_{0,k}(A;J)$ 表示其中 $u$ 单纯的子集。对一般的 $J$，它是
维数为
\[
  \dim\mathcal M^*_{0,k}(A;J)
  =2n+2c_1(A)+2k-6=:\mu(A,k)
\]
的流形。此外，有显然的\BookRecordIndex{T0283}求值映射
\[
  \widetilde{\mathcal M}_{0,k}(A;J)\longrightarrow M^k,
  \qquad (u,z_1,\dots,z_k)\longmapsto\bigl(u(z_1),\dots,u(z_k)\bigr),
\]
它下降为
\[
  \operatorname{ev}:\mathcal M_{0,k}(A;J)\longrightarrow M^k.
\]
如果这个映射表示一个循环，就可用它与 $M^k$ 中循环的交数定义\BookRecordIndex{T0143}格罗莫夫--威滕不变量。

请注意，尽管 $\mathcal M_{0,k}(A;J)$ 由无参数映射组成，但当 $k\ge2$ 时它绝不会紧致，
因为按定义各\BookRecordIndex{T0066}标记点彼此不同。尽管如此，我们希望用\BookRecordIndex{T0283}求值映射
\[
  \operatorname{ev}:\mathcal M^*_{0,k}(A;J)\longrightarrow M^k
\]
所表示的 $d$ 维同调类来定义\BookRecordIndex{T0143}格罗莫夫--威滕不变量，其中 $d=\mu(A,k)$。因此，即使已经
知道 $\mathcal M_{0,k}(A;J)$ 是流形，仍须证明它表示一个循环。关键在于理解其紧化。
Kontsevich 认识到，结构最好的 $\mathcal M_{0,k}(A;J)$ 紧化并不是格罗莫夫考虑的\BookRecordIndex{T0181}尖点
曲线空间，而是\BookRecordIndex{T0340}稳定映射空间 $\overline{\mathcal M}_{0,k}(A;J)$；见定义 5.1.1。
第 5 章使用这种语言表述并证明格罗莫夫紧致性定理。

即使很好地理解了紧化 $\overline{\mathcal M}_{0,k}(A;J)$，对一般流形 $(M,\omega)$，
仍很难说明 $\operatorname{ev}$ 在何种意义下表示一个循环，因为低\BookRecordIndex{T0072}层中的\BookRecordIndex{T0113}多重覆盖曲线可能
使这些\BookRecordIndex{T0072}层的维数过高。
\footnote{这一问题如今已经在一般流形上得到解决。通过精细的整体扰动，可以定义一个
$d$ 维\BookRecordIndex{T0370}虚模循环，应把它看作 $\mathcal M^*_{0,k}(A;J)$ 的基本类。详细参考文献见第 6 章
导言。这超出本书范围。}
不过，如果 $(M,\omega)$ \BookRecordIndex{T0058}半正，则第 6 章将证明\BookRecordIndex{T0283}求值映射表示一个维数为
$d:=\mu(A,k)$ 的\BookRecordIndex{T0337}伪循环。直观地说，这表示可以通过添加一个维数至多为 $d-2$ 的集合来
紧化它的像。
\footnote{光滑流形 $X$ 的子集 $B$，若包含在某个定义于 $m$ 维流形 $W$ 上的光滑映射
$g:W\to X$ 的像中，就称其维数至多为 $m$。}
这样，它的边界维数至多为 $d-2$，从同调观点看不见这个边界。更正式地，有如下定义。

\begin{nameddefinition}
流形 $X$ 中的一个 $d$ 维\BookRecordIndex{T0337}伪循环，是定义在定向 $d$ 维流形 $V$ 上的光滑映射
\[
  f:V\to X,
\]
使得像 $f(V)$ 的闭包紧致，并且其极限集 $\Omega_f$ 的维数不超过 $d-2$。这里
\[
  \Omega_f=
  \bigcap_{\substack{K\subset V\\K\text{ 紧}}}
  \overline{f(V\setminus K)}
\]
是所有序列 $f(x_\nu)$ 的极限点组成的集合，其中 $x_\nu$ 在 $V$ 中没有收敛子列。
\end{nameddefinition}

第 6.5 节进一步研究这个定义。下面的定理在第 6.6 节证明。回忆，$\mathcal J$ 表示所有
被 $\omega$ \BookRecordIndex{T0372}驯服的 $J$ 组成的空间。

\begin{namedtheorem}[C]
设 $(M,\omega)$ 是闭\BookRecordIndex{T0058}半正\BookRecordIndex{T0358}辛流形。则存在一个\BookRecordIndex{T0300}剩余集
$\mathcal J_{\mathrm{reg}}(M,\omega)\subset\mathcal J_\tau(M,\omega)$，具有下述性质。
若 $A\in H_2(M;\mathbb Z)$ 满足 $c_1(A)>0$，并且
$J\in\mathcal J_{\mathrm{reg}}(M,\omega)$，则\BookRecordIndex{T0283}求值映射
\[
  \operatorname{ev}:\mathcal M^*_{0,k}(A;J)\longrightarrow M^k
\]
是维数为 $\mu(A,k)$ 的\BookRecordIndex{T0337}伪循环。此外，它的同调类与
$J\in\mathcal J_{\mathrm{reg}}(M,\omega)$ 的选择无关。
\end{namedtheorem}

证明的梗概见第 6 章导言，其中要用到第 6.1--6.3 节建立的\BookRecordIndex{T0283}求值映射\BookRecordIndex{T0166}横截性结果。尤其建议
读者仔细研究第 6.2 节如何用\BookRecordIndex{T0065}边求值映射的\BookRecordIndex{T0166}横截性精确定义
$\mathcal J_{\mathrm{reg}}(M,\omega)$；没有这个精确表述，就无法充分理解定理 C。

定理 C 有一个重要变体。考虑 $\widetilde{\mathcal M}^*_{0,k}(A;J)$ 的子空间
$\widetilde{\mathcal M}^*_{0,k}(A;\mathbf w,J)$；其中元组
$(u,z_1,\dots,z_k)$ 满足：与指标 $i\in I$ 对应的\BookRecordIndex{T0066}标记点 $z_i$ 被固定为 $w_i$。
（这里假设 $\#I\ge3$，并记 $\mathbf w:=\{w_i\}_{i\in I}$。）为了证明所得到的映射
\[
  \operatorname{ev}:\mathcal M^*_{0,k}(A;\mathbf w,J)\longrightarrow M^k
\]
是\BookRecordIndex{T0337}伪循环（这一次其维数为 $\mu(A,k)-2(\#I-3)$），必须引入更广的一类\BookRecordIndex{T0209}柯西--黎曼方程
扰动，考察 $J=\{J_z\}$ 随 $z\in S^2$ 变化的方程之解。从几何上说，这等价于考察适当
映射 $u:S^2\to M$ 的图像 $z\mapsto(z,u(z))$。当 $c_1(A)=0$ 时，为了证明上面的定理 C，
也需要类似方法。要把定理 C 推广到这些情形，就必须在稍广的框架中重做整套理论。细节见
第 6.7 节。
